% !TeX spellcheck = en_GB
\documentclass[a4paper,12pt]{book}
\usepackage{polyglossia} % Support for multiple languages
%\setmainlanguage{english} % Default language
\setotherlanguage{arabic} % Secondary language
\usepackage{fontspec}
\setmainfont{Times New Roman}
\newfontfamily\arabicfont[Script=Arabic]{Simplified Arabic} % Arabic font
\usepackage{graphicx}
\usepackage[absolute,overlay]{textpos}
\usepackage{setspace}
\usepackage{titlesec}
\usepackage{geometry}
\geometry{
	right=3.5cm,
	left=2.5cm,
	top=2.5cm,
	bottom=2.5cm,
}
\usepackage{tocloft}

\renewcommand{\cftchapdotsep}{\cftdotsep} % Add dots for chapters
\renewcommand{\cftsecdotsep}{\cftdotsep} % Add dots for sections
\renewcommand{\cftsubsecdotsep}{\cftdotsep} 

\usepackage{fancyhdr}
\pagestyle{fancy}
\fancyhf{}
\fancyfoot[L]{\thepage} 

\renewcommand{\headrulewidth}{0.5pt} % Remove the header rule
\renewcommand{\footrulewidth}{0.5pt} 

\usepackage{lmodern}
\usepackage{xcolor}

% Line spacing
\onehalfspacing

% Section styles
\titleformat{\section}
{\normalfont\fontsize{18}{18}\bfseries}{\thesection}{1em}{}
\titleformat{\subsection}
{\normalfont\fontsize{16}{16}\bfseries}{\thesubsection}{1em}{}
\titleformat{\subsubsection}
{\normalfont\fontsize{14}{14}\bfseries\itshape}{\thesubsubsection}{1em}{}

%\graphicspath{{../../../../Assets/img/}}

\raggedbottom % Avoid LaTeX adding vertical space

\let\cleardoublepage\clearpage

\begin{document}
	\color{black}
	
%===========================================================
%Cover Page (Arabic)
%===========================================================
	\begin{titlepage}
		\begin{flushright}
			\fontsize{24}{24}\selectfont\textbf{\textarabic{جامعة طرابلس}}\\[10pt]
		\end{flushright}
		\hfill    
		\begin{textblock*}{3cm}(1cm,1cm)
			\includegraphics[width=3cm]{UOT_logo_transp.png}
		\end{textblock*}
		\vfill
		\begin{center}
			\fontsize{24}{24}\selectfont\textbf{\textarabic{عنوان مشروع التخرج}}\\[20pt]
			\fontsize{14}{14}\selectfont{\textarabic{سند العروسي \hspace*{4em} 2181801442} }\\[10pt]
			\vfill
			\fontsize{20}{20}\selectfont\textbf{\textarabic{بكالوريوس العلوم في تقنية المعلومات}}\\[10pt]
			\fontsize{14}{14}\selectfont{\textarabic{الفصل الدراسي}}\\
		\end{center}
	\end{titlepage}
	
%===========================================================
%Title Page
%===========================================================
	\begin{titlepage}
		\begin{flushleft}
			\hfill
			\fontsize{24}{24}\selectfont\textbf{\textarabic{جامعة طرابلس}}\\[10pt]
			\hfill
			\fontsize{22}{22}\selectfont\textbf{\textarabic{كلية تقنية المعلومات}}
		\end{flushleft}
		\hfill	
		\begin{textblock*}{3cm}(1cm,1cm)
			\includegraphics[width=4cm, height=4cm]{Faculty_logo_transp.png}
		\end{textblock*}
		\vfill
		\begin{center}
			\fontsize{24}{24}\selectfont\textbf{\textarabic{عنوان المشروع}}\\[20pt]			
			\fontsize{14}{14}\selectfont{\textarabic{مشروع مقدم}}\\
			\fontsize{14}{14}\selectfont{\textarabic{لقسم هندسة البرمجيات}}\\
			\fontsize{14}{14}\selectfont{\textarabic{لإستيفاء المتطلبات الجزئية لنيل درجة}}\\ \fontsize{14}{14}\selectfont{\textarabic{البكالوريوس في علوم}}\\
			\fontsize{14}{14}\selectfont{\textarabic{تقنية المعلومات}}\\[10pt]
			\vfill
			\fontsize{14}{14}\selectfont{\textarabic{من قبل الطالب:}}\\[10pt]
			\fontsize{14}{14}\selectfont{\textarabic{سند العروسي \hspace*{4em} 2181801442} }\\[10pt]
			\vfill
			\fontsize{14}{14}\selectfont{\textarabic{بإشراف: }}\\
			\fontsize{14}{14}\selectfont{\textarabic{الدكتور/الأستاذ/المحاضر فلان}}
			\vfill
			%\fontsize{20}{20}\selectfont\textbf{Bachelor of Information Technology}\\[10pt]
			\fontsize{14}{14}\selectfont{\textarabic{الفصل الدراسي}}\\
		\end{center}
	\end{titlepage}
	
	

	
	\begin{otherlanguage}{arabic}
%===========================================================
%Dedication الإهداء
%===========================================================
	\chapter*{الإهداء}
	\addcontentsline{toc}{chapter}{الإهداء}
	نص الإهداء هنا.

%===========================================================
%Akcnowledgments الشكر و التقدير
%===========================================================
	\chapter*{الشكر والتقدير}
	\addcontentsline{toc}{chapter}{الشكر والتقدير}
	نص الشكر والتقدير هنا.

%===========================================================
%AR Abstract الملخص بالعربي
%===========================================================
	\chapter*{الملخص بالعربية}
	\addcontentsline{toc}{chapter}{الملخص بالعربية}
	نص الملخص بالعربية هنا.
\end{otherlanguage}


%===========================================================
%EN Abstract الملخص بالانجليزي
%===========================================================
	\chapter*{Abstract}
	\begin{otherlanguage}{arabic}
	\addcontentsline{toc}{chapter}{الملخص English in Abstract} %لازم هكي باش يقعد يمين
\end{otherlanguage}
	English abstract text here.

\newpage
\begin{otherlanguage}{arabic}
	% Table of Contents
	\renewcommand{\contentsname}{المحتويات} % Change "Contents" to Arabic
	
	\tableofcontents
	
\end{otherlanguage}

%===========================================================
%Introduction المقدمة
%===========================================================
	\begin{otherlanguage}{arabic}
	

%===========================================================
%Start writing your grad documentation from here التوثيق يبدأ من هنا
%===========================================================
	\chapter*{المقدمة}
	\addcontentsline{toc}{chapter}{المقدمة}
	نص المقدمة هنا.	الباب الأول
	
	\section{القسم الرئيسي}
	هذا هو القسم الرئيسي من الباب الأول.
	\paragraph{مرحبا جميعاً انا باراقراف}
	\subparagraph{وانا باراقراف فرعي}
	
	
	\subsection{القسم الفرعي ١}
	هذا هو القسم الفرعي الأول من القسم الرئيسي الأول من الباب الأول.
	
	\subsubsection{القسم الفرعي ٢}
	هذا هو القسم الفرعي الثاني من القسم الفرعي.
	
	\subsubsection{القسم الفرعي ٣}
	هذا هو القسم الفرعي الثالث من القسم الفرعي.
	
	
	%chapters
	%sections
	
%===========================================================
%EoDocumentation نهاية التوثيق
%===========================================================
	\chapter*{الخاتمة}
	\addcontentsline{toc}{chapter}{الخاتمة}
	نص الخاتمة هنا.
	
%===========================================================
%References المراجع
%===========================================================
	\chapter*{المراجع}
	\addcontentsline{toc}{chapter}{المراجع}
	صفحة المراجع
	
	\begin{thebibliography}{9}
		\addcontentsline{toc}{chapter}{المصادر}
		\bibitem{ref1} مرجع 1
		\bibitem{ref1} مرجع 2
		\bibitem{ref1} مرجع 3
	\end{thebibliography}
	 
%===========================================================
%Appendices الملاحق
%===========================================================
	\chapter*{الملاحق}
	\addcontentsline{toc}{chapter}{الملاحق}
	نص الملاحق هنا.
	
	
	% Arabic Section
	
	\end{otherlanguage}
	
\end{document}
